\documentclass{article}

\usepackage{amsmath,amssymb}
\usepackage{xurl}
\usepackage[hidelinks]{hyperref}

% Shared commands for notes in docs/paper-gaps/.

\DeclareUnicodeCharacter{2080}{\ifmmode{}_{0}\else\textsubscript{0}\fi}
\DeclareUnicodeCharacter{2081}{\ifmmode{}_{1}\else\textsubscript{1}\fi}
\DeclareUnicodeCharacter{2082}{\ifmmode{}_{2}\else\textsubscript{2}\fi}
\DeclareUnicodeCharacter{2099}{\ifmmode{}_{n}\else\textsubscript{n}\fi}

\newcommand{\Lean}{\textsc{Lean}}
\ExplSyntaxOn
\NewDocumentCommand{\leanid}{m}{
  \ifmmode
    \textsf{\detokenize{#1}}
  \else
    \group_begin:
    \urlstyle{sf}
    \tl_set:Nn \l_tmpa_tl {#1}
    \tl_replace_all:Nnn \l_tmpa_tl {\_} {_}
    \exp_args:NV \path \l_tmpa_tl
    \group_end:
  \fi
}
\ExplSyntaxOff
\providecommand{\tr}{\operatorname{tr}}
\newcommand{\tnleanIssueBase}{https://github.com/LionSR/TNLean/issues/}
\newcommand{\tnleanPrBase}{https://github.com/LionSR/TNLean/pull/}
\newcommand{\tnleanDocBase}{https://sirui-lu.com/TNLean/docs/}
\newcommand{\ghissue}[1]{\href{\tnleanIssueBase#1}{GitHub issue \##1}}
\newcommand{\ghpr}[1]{\href{\tnleanPrBase#1}{PR \##1}}
\newcommand{\doclink}[1]{\href{\tnleanDocBase#1.html}{\path{#1}}}

% Dirac notation and the identity operator, as in blueprint/src/macros/common.tex.
% \providecommand keeps note-local definitions authoritative.
\usepackage{dsfont}
\providecommand{\ket}[1]{|#1\rangle}
\providecommand{\bra}[1]{\langle #1|}
\providecommand{\braket}[2]{\langle #1 | #2 \rangle}
\providecommand{\ketbra}[2]{|#1\rangle\!\langle #2|}
\providecommand{\Id}{\mathds{1}}
\providecommand{\one}{\mathds{1}}
\providecommand{\C}{\mathbb{C}}
\providecommand{\MN}[1]{M_{#1}(\C)}
\providecommand{\MD}{M_D(\C)}

% External referencing for paper sources.  The build helper writes sanitized
% auxiliary files under build/paper-aux/ and excludes duplicated source labels.
\usepackage{xr}

% Import a generated paper auxiliary file with a caller-supplied label prefix.
% The candidate paths support builds from docs/paper-gaps/, the repository root,
% and build/paper-aux/ itself.
\newcommand{\importpaperaux}[2]{%
  \IfFileExists{../../build/paper-aux/#2.aux}{%
    \externaldocument[#1]{../../build/paper-aux/#2}%
  }{%
    \IfFileExists{build/paper-aux/#2.aux}{%
      \externaldocument[#1]{build/paper-aux/#2}%
    }{%
      \IfFileExists{#2.aux}{%
        \externaldocument[#1]{#2}%
      }{}%
    }%
  }%
}

% General external reference macros
\newcommand{\paperref}[2]{\ref{#1-#2}}
\newcommand{\papereqref}[2]{(\ref{#1-#2})}

% CPSV16 source (1606.00608 / MPDO-22-12-17-2.tex) external document and shortcuts
\importpaperaux{cpsv16-}{1606.00608/MPDO-22-12-17-2-tnlean}
\newcommand{\cpsvref}[1]{\paperref{cpsv16}{#1}}
\newcommand{\cpsveqref}[1]{\papereqref{cpsv16}{#1}}

% Verdict marker: \gapnote{<kind>}{<status>}. Kind is the policy
% classification (clarification, local-correction, scope-restriction,
% unfaithful, false-source, open-gap); status is open, wip (elimination
% underway), resolved, or historical. Machine-read by the site index;
% prints a verdict line.
\newcommand{\gapnote}[2]{\par\noindent\textbf{Verdict:} #1 (#2).\par}


\title{Wolf Theorem~6.14: Fixed-Point Projection Gap}
\author{}
\date{2026-07-18}

\begin{document}
\maketitle

\section{The completed classification theorem}

Wolf's Proposition~1.5 assumes a positive linear retraction onto a unital
finite-dimensional $*$-subalgebra and classifies that retraction
\cite[Proposition~1.5 and Equation~(1.40)]{Wolf2012Quantum}.  The local source
is \path{Notes/WolfNoteTexSource/ch01_deconstructing_quantum.tex},
lines~536--570.

The theorem
\leanid{StarSubalgebra.exists_block_densities_of_positive_retraction} proves
the full proposition.  It assumes exactly that the given linear map is
positive, that its image is contained in the subalgebra, and that it fixes
every member of the subalgebra.  It obtains the unitary direct-sum
representation, removes off-diagonal inputs by positivity, and gives a
positive semidefinite trace-one matrix for each summand.  No assumption of
complete positivity, trace preservation, or a Schwarz inequality is added.

\section{The remaining fixed-point theorem}

Let $T:M_D(\mathbb C)\to M_D(\mathbb C)$ be positive and trace preserving,
and suppose that its trace-pairing adjoint satisfies the Schwarz inequality.
Wolf's Theorem~6.14 asserts that there are a unitary $U$, a decomposition
\begin{align}
    \mathbb C^D
      &= \mathcal H_0\oplus
        \bigoplus_{k=1}^K\mathbb C^{d_k}\otimes\mathbb C^{m_k},
\end{align}
and positive definite density matrices $\rho_k\in M_{m_k}(\mathbb C)$ such
that
\begin{align}
    \operatorname{Fix}(T)
      &= U\left(0\oplus\bigoplus_{k=1}^K
          M_{d_k}(\mathbb C)\otimes\rho_k\right)U^*.
    \tag{Wolf, Theorem 6.14 and Equation (6.63)}
\end{align}
Here $\mathcal H_0$ is a genuine zero summand: every fixed point vanishes on
this subspace.  This is the statement in
\cite[Theorem~6.14]{Wolf2012Quantum}; the local source is
\path{Notes/WolfNotePDF/wolf_ch6_sections.txt}, lines~896--905.

The completed formal theorem classifies a given positive retraction onto a
unital $*$-subalgebra.  Three further arguments are needed to obtain the
fixed-point statement above.  First, the hypotheses on $T$ must be used to
construct a positive retraction onto $\operatorname{Fix}(T)$, for example by
Ces\`aro means, and to identify the full-support part of this fixed-point
space with a unital $*$-algebra.  Second, the positive semidefinite density
matrices supplied by Proposition~1.5 must be restricted to their supports;
this produces the positive definite matrices in Theorem~6.14.  Third, the
discarded kernels must be assembled with the subspace on which every fixed
point vanishes, thereby recovering the zero summand $\mathcal H_0$.

Thus Proposition~1.5 is fully formalized, while Theorem~6.14 remains an open
gap.  The missing work is not another classification of a given retraction:
it is the construction of the fixed-point retraction, the full-support
reduction, and the recovery of the zero block from the kernels.

\bibliographystyle{alpha}
\bibliography{references}

\end{document}
